\section*{A097: rowwise minimum cut for the complete singleton KL layer}

All logarithms are base two.  Let $f:\{-1,1\}^n\to\{-1,1\}$ be a
genuine Boolean function.  Put $p(S)=\widehat f(S)^2$ and
$\pi(A)=\sum_{S\supseteq A}p(S)$.  In the exact signed restriction identity
write
\[
 G=D_{\ge2}+K^{(0)}+K^{(1)},\qquad
 D_{\ge2}=\sum_i(B_i-\Gamma_i)+\sum_i\sum_{|A|\ge2}K_{i,A}.       \tag{1}
\]
Thus the signed term and every KL row of order at least two remain intact.

For a singleton $A=\{j\}$ define
\[
 r_j=\sum_{i\ne j}p(ij),\qquad
 s_j=\sum_{\substack{|S|\ge3\\j\in S}}p(S),                    \tag{2}
\]
\[
 V_j=\sum_{i\ne j}(p(j)+p(ij))
 \log_2\frac{\pi(j)}{p(j)+p(ij)}.                               \tag{3}
\]
Zero-mass summands are omitted.  Finally put
\[
 M_1=\sum_i\min\{p(\varnothing),p(i)\},\qquad c_0={2\over\ln2}. \tag{4}
\]

\subsection*{The rowwise minimum theorem}

For every $j$ the complete row with lower set $\{j\}$ obeys both
\[
 \sum_{i\ne j}K_{i,\{j\}}
 \le {1\over2}V_j+{1\over\ln2}r_j                               \tag{5}
\]
and
\[
 \sum_{i\ne j}K_{i,\{j\}}
 \le {9\over\ln2}(2r_j+s_j).                                   \tag{6}
\]
Consequently
\[
 \boxed{
 K^{(1)}\le\sum_j\min\left\{
 {1\over2}V_j+{r_j\over\ln2},
 {9\over\ln2}(2r_j+s_j)\right\}.}                              \tag{7}
\]
The minimum is taken before summing over $j$; it is therefore at least as
strong as the minimum of the two previously available global payments.

Equation (5) is the genuine-child singleton estimate.  For completeness,
fix $i\ne j$, write $a=\widehat f(j)$, $b=\widehat f(ij)$,
$m=a^2+b^2$, $M=\pi(j)$, and retain the signed upper correlation $C$.
Boolean granularity in the two physical children gives
\[
 K_{i,\{j\}}\le {1\over2}m\log_2{M\over m}
                   +{1\over\ln2}p(ij).                          \tag{8}
\]
Summing in $i$ proves (5).

We include the Boolean-specific proof of (6), because retaining its sharper
level weights is the new assembly step.  Form the physical derivative
\[
 g=D_jf\in\{-1,0,1\}.
\]
For a coordinate $i\ne j$ set
\[
 a=\widehat g(\varnothing),\quad b=\widehat g(i),\quad
 M=\mathbb E g^2,\quad C={1\over2}\widehat{g^2}(i),\quad m=a^2+b^2.
\]
The identity $|g|=g^2$ gives the signed branch constraints
\[
 |a\pm b|\le M\pm2C.                                             \tag{9}
\]
With $A=a/M$, $B=b/M$, $\gamma=2C/M$, $S=A^2+B^2$, the row is
$mD_2(P\|Q)$, where
\[
 P_\pm={(A\pm B)^2\over2S},\qquad Q_\pm={1\pm\gamma\over2}.
\]
If $|\gamma|\le1/3$, $D\le\chi^2$ and
\[
 S D(P\|Q)\le{(2AB-\gamma S)^2\over S(1-\gamma^2)}
 \le9(B^2+\gamma^2).                                             \tag{10}
\]
If $|\gamma|>1/3$, (9) yields $P_\pm/Q_\pm\le2/S$, so
\[
 S D(P\|Q)\le S\log(2/S)\le2/e<9\gamma^2.                       \tag{11}
\]
Converting natural logarithms to base two and multiplying by $M^2$ gives
\[
 K_{i,\{j\}}\le {9\over\ln2}(b^2+4C^2).                        \tag{12}
\]
After summing in $i$, Parseval and $|\widehat g(\varnothing)|\le M$ give
\[
 \sum_{i\ne j}K_{i,\{j\}}
 \le {9\over\ln2}\bigl(2W_1(g)+W_{\ge2}(g)\bigr)
 ={9\over\ln2}(2r_j+s_j),                                      \tag{13}
\]
which is (6).  Notice that no unsigned replacement of $C$ occurs.

Summing (6) alone retains the sharper global weights
\[
 \boxed{K^{(1)}\le {1\over\ln2}
 \left(36W_2+9\sum_{d\ge3}dW_d\right).}                         \tag{14}
\]
This improves the coarser $36(W_2+\sum_{d\ge3}dW_d)/\ln2$ statement.

\subsection*{New universal row cut and exact unpaid leaf}

The sharp root theorem $K^{(0)}\le c_0M_1$, (1), and (7) give
\[
 \boxed{
 G\le D_{\ge2}+c_0M_1+
 \sum_j\min\left\{{V_j\over2}+{r_j\over\ln2},
 {9\over\ln2}(2r_j+s_j)\right\}.}                              \tag{15}
\]
This is universal, dimension independent, dummy invariant, and preserves
all signed, normalization, and higher-row terms.

Freeze $R=G-4I-3U/4$.  Tensor additivity, the degree-at-most-two theorem,
and dimension-minimal descent show that it is sufficient to prove
\[
 \boxed{
 D_{\ge2}+c_0M_1+
 \sum_j\min\left\{{V_j\over2}+{r_j\over\ln2},
 {9\over\ln2}(2r_j+s_j)\right\}
 \le4I+{3\over4}U}                                               \tag{16}
\]
only for tensor-prime genuine Boolean functions with $R>0$ and Fourier
degree at least three.  Combining (15) and (16) gives the frozen GU pair.
The audited restriction induction then gives $U\le16I$, and
$H\le2I+2U$ gives $H\le34I$.

Equation (16), not universal GU, is the remaining obligation.  Relative to
the previous high-degree-reserve leaf it replaces the global $V_1/2$ charge
by a rowwise minimum.  In particular, the one-point-surgery family for
which $V_1/2$ grows linearly while $Q_{\ge3}$ stays bounded no longer creates
an unbounded singleton charge: (14) depends only on genuine level-two and
higher spectral energy.  The signed $D_{\ge2}$ term remains unpaid and must
not be separated from the final estimate.

For that exact surgery spectrum, the certified lower bound on the former
global singleton budget exceeds (14) already at defect parameter $d=50$:
$13.2213475204>12.9842553680$.  The former bound then grows linearly in
$d$, whereas (14) converges to $9/\ln2$.  Hence the new cut is strictly
stronger on the precise family that invalidated a separate $V_1$ payment.

\subsection*{Verified conditional reduction to FEI}

For clarity we record the induction used above.  Suppose universally
$G\le aI+bU$ with $b<1$, and put $C=a/(1-b)$.  Constants have
$I=U=0$.  For a nonconstant $f$, since
$\sum_i\Delta_iU=G$ and $\sum_i\operatorname{Inf}_i=I$, some relevant
coordinate $i$ satisfies
\[
 \Delta_iU\le {aI+bU\over I}\operatorname{Inf}_i.                \tag{17}
\]
Induction on the two physical children gives
\[
 {U(f_{i,+})+U(f_{i,-})\over2}
 \le C\bigl(I-\operatorname{Inf}_i\bigr).                        \tag{18}
\]
Writing $r=U/I$ and $t=\operatorname{Inf}_i/I>0$, (17)--(18) imply
\[
 r\le C(1-t)+(a+br)t=C+bt(r-C).
\]
Hence $(1-bt)(r-C)\le0$, and $r\le C$.  Thus $U\le aI/(1-b)$.
The separately proved source inequality
$H\le2(I-\operatorname{Var}f)+2U\le2I+2U$ then gives
\[
 H\le\left(2+{2a\over1-b}\right)I.                              \tag{19}
\]
For $(a,b)=(4,3/4)$ this is $H\le34I$.
